\documentclass[12pt,a4paper]{article}

% encoding / language
\usepackage[utf8]{inputenc}
\usepackage[T1]{fontenc}
\usepackage[brazilian]{babel}

% misc
\usepackage{fmtcount}
\usepackage{amsmath,amssymb,amsthm}
\usepackage{enumitem}
\usepackage[normalem]{ulem}
\usepackage{fancyvrb}
\usepackage{xcolor}

% citations
\usepackage[style=abnt]{biblatex}
\addbibresource{ref.bib}

% spacing
\frenchspacing
\setlength{\parindent}{1.25cm}
\usepackage{setspace}
\onehalfspacing
\usepackage[top=3cm,bottom=2cm,left=3cm,right=2cm]{geometry}

% font
\usepackage[scaled]{helvet}
\renewcommand{\familydefault}{\sfdefault}
\usepackage{sfmath}

\newenvironment{question}[1]{
    \par\medskip\noindent
    \textbf{Questão #1} \dotfill
    \par\nopagebreak\medskip
}{
    \medskip
}

\newcommand{\listtitle}[1]{
    \par\medskip\noindent
    {\rmfamily\small\textbf{\uline{#1:}}}
    \par\nopagebreak
}

\renewcommand{\theFancyVerbLine}{\ttfamily\color{gray}\arabic{FancyVerbLine}}

\title{Resolução de questões objetivas de Complexidade de Algoritmos do Enade de Ciência da Computação}
\date{\today}

\begin{document}
\maketitle

\ordinalnum{1}[f] atividade da disciplina optativa ENADE-2026.2, do Professor Ricardo Azevedo Moreira da Silva.
As questões objetivas (múltipla escolha) foram selecionadas das últimas três provas do Exame Nacional de Desempenho dos Estudantes (Enade)
e envolvem o assunto de Complexidade de Algoritmos.
A seção \textbf{Fundamentos} contém todo o conteúdo necessário e devidamente resumido para um estudante de Ciência da Computação resolver as questões.
A seção \textbf{Questões} contém as questões resolvidas e justificadas.
Essas seções citam livros presentes na seção \textbf{Referências}.

\section{Fundamentos}
Segundo \textcite[35]{clrs2022}, blá blá blá...
Isso resulta em uma complexidade de tempo de $O(n)$ \parencite{clrs2022}.
\dots

\section{Questões}
\begin{question}{2014-12}
Analise o custo computacional dos algoritmos a seguir, que calculam o valor de um polinômio de grau $n$, da forma:
$P(x) = a_n x^n + a_{n-1} x^{n-1} + \dots + a_1 x + a_0$,
onde os coeficientes são números de ponto flutuante armazenados no vetor $a[0..n]$, e o valor de $n$ é maior que zero.
Todos os coeficientes podem assumir qualquer valor, exceto o coeficiente $a_n$, que é diferente de zero.

\noindent\begin{minipage}{\linewidth}
\listtitle{Algoritmo 1}
\begin{Verbatim}[numbers=left]
soma = a[0]
Repita para i = 1 até n
    Se a[i] != 0.0 então
        potência = x
        Repita para j = 2 até i
            potência = potência * x
        Fim repita
        soma = soma + a[i] * potencia
    Fim se
Fim repita
Imprima(soma)
\end{Verbatim}

\listtitle{Algoritmo 2}
\begin{Verbatim}[numbers=left]
soma = a[n]
Repita para i = n-1 até 0 passo -1
    soma = soma * x + a[i]
Fim repita
Imprima(soma)
\end{Verbatim}
\medskip
\end{minipage}

\noindent
Com base nos algoritmos 1 e 2, avalie as asserções a seguir e a relação proposta entre elas.

\begin{enumerate}[label=\Roman*.]
    \item Os algoritmos possuem a mesma complexidade assintótica.
\end{enumerate}

\begin{center}
    \textbf{PORQUE}
\end{center}

\begin{enumerate}[label=\Roman*., start=2]
    \item Para o melhor caso, ambos os algoritmos possuem complexidade $O(n)$.
\end{enumerate}

\medskip\noindent
A respeito dessas asserções, assinale a opção correta.

\begin{spacing}{0.9}
\begin{small}
\begin{enumerate}[label=\Alph*)]
    \item As asserções I e II são proposições verdadeiras, e a II é uma justificativa correta da I.
    \item As asserções I e II são proposições verdadeiras, mas a II não é uma justificativa correta da I.
    \item A asserção I é uma proposição verdadeira, e a II é uma proposição falsa.
    \item A asserção I é uma proposição falsa, e a II é uma proposição verdadeira.
    \item As asserções I e II são proposições falsas.
\end{enumerate}
\end{small}
\end{spacing}
\end{question}

\begin{question}{2014-16}
\dots
\end{question}

% 2014-22
% 2014-28
% 2014-30
% 2017-25
% 2017-33
% 2021-20
% 2021-32

\printbibliography
\end{document}
